
% Shoggoth 22 Aug 26 


\subsection{The Congregation and Lay Activity}

Editorial article in “Folkebladet” for April 10, 1907.

In the previous number of “Folkebladet” we had to write something about the manner in which the large Association Committee dealt with “Lay Activity.” It is unpleasant enough to have to take up such misleading representations of a precious cause; and we therefore owe it to ourselves this time to speak more positively about the matter.

When, then, we wish to speak about the congregation and lay activity, it is in order to say that it is to be expected that it will make a difference also for the cause of lay activity where there is a congregation.

By a congregation we mean a gathering of working Christians who use the means of grace and the gifts of grace for their own salvation and blessedness and that of their fellow human beings. Where such a gathering of active Christians is found, the question will not be one of binding and controlling lay activity and other gifts of grace, but much rather of giving them room and freedom, opportunity and encouragement, so that they may attain their greatest flourishing and most effective use.

Here, then, a question arises which there was no occasion to discuss in the previous article, namely: who has this gift of speaking in an assembly, which is, after all, one of the conditions of lay activity? And who is to decide whether such a gift is present? And how is such a gift best to be discovered and developed?

Here lies a difficulty which all have felt who have thought with interest and love about lay activity and its flourishing. Not everyone has received the gift merely because he imagines that he has it. Nor does anyone receive the gift by being called by a congregation to speak in an assembly. It is evident, indeed, in the case of a number of pastors that they themselves believe they have the gift, and they also have the call of the congregation, but nevertheless the gift of grace does not appear to be present.

With regard to this matter we have heard that Hauge’s friends had the custom that some of the older, tested Christians and preachers of God’s Word exercised the work of testing spirits and gifts, and that new and younger forces had to submit to a fairly long period of testing before they received an independent place in the ranks of the witnesses. And a similar procedure should surely be sought also among us. We say “similar,” for precisely the same manner will probably be difficult to attain.

It is clear that neither chairs nor examinations can produce a lay preacher; but neither does anyone become a lay preacher simply by giving up other work without further ado and making it his livelihood to travel about and preach. A traveling evangelist or revivalist who has this work as his sole occupation and livelihood is not a lay preacher in the older sense of the word. We understand also that those who make the work of preaching their only occupation are preferably called emissaries, both in Norway and here among us.

Without doubt the work of an evangelist has its great significance. But it is not Haugean lay activity. We do not depreciate the one work merely because we write in commendation of the other.

A lay preacher is a preacher who has a message from God to bring to his fellow human beings, and who cannot remain silent concerning that message which God has given him to bear forth without unfaithfulness; and who therefore, both in conversation with individuals and in larger or smaller assemblies of various kinds, brings forth this message. But he continues to be a layman while he thus preaches, and does not derive his livelihood from this work of preaching.

They are driven to speak and cannot refrain from it; they begin wherever they have opportunity, in their home and in their home district; they travel about for shorter or longer distances, and neither congregational resolutions nor other prohibitions bar their way.

Naturally, they are dependent upon what trust they gain, especially among other Christians. To give them a congregational call is not necessary and therefore not desirable either; for then they could regulate neither their time nor their preaching according to their ability and gift, but would have to try to please those in whose pay they were. But this does not exclude the congregation from encouraging and urging them on in such a manner as it itself finds right and useful.

But their gift is tested through the use of the gift and is proved by the reception they gain. There are honorable and upright ways by which their influence may gradually be extended into wider circles, if they attain that glorious thing of gaining the trust of Christians. But to procure for oneself a whole number of testimonials, which perhaps are given half unwillingly because no one wishes to offend them by refusing them a testimonial, is often equivalent to a lack of real trust.

The true lay preacher is independent of monetary gifts. He has his temporal work and his livelihood apart from his Christian witness. Therefore his authority is greater, and his witness is often more effective among many people.

Such men stand as free and independent as is at all compatible with the Christian estate and brotherhood. And in this lies, in no small degree, their value to the congregation. They belong to the congregation, and they belong to the Lord; they are driven by the Lord’s Spirit to bear witness to those within the congregation and to those outside the congregation. The congregation cannot, without overreach, command them to be silent when it must be acknowledged that God has given them a message and a gift; neither can they demand of the congregation that it support them when it has not called them. But where there is a free and living congregation, it may be expected both that there will be much such lay activity and that broad room and free opportunity will be opened for it. For although not many in each place have great gifts in this direction, many Christians ought nevertheless to have spirit and power to bear such witness as might be for the edification of those who hear it.

We are altogether convinced that living congregations will do well to heed the apostolic rule in 1 Corinthians 14:

“What then is to be done, brothers? When you come together, each of you has a psalm, he has a teaching, he has a word, he has an interpretation; let all things be done for edification!”

And it is surely not too much to say that if, according to this apostolic rule, all Christians become the Savior’s witnesses also “in an assembly,” lay activity too will flourish better. For, first, the gifts will come into use in the proper Christian manner; and second, a lay preacher will not come to stand so alone as a “lone tree” on the prairie, which is always a dangerous and sorrowful position.

Let a little life come into the work, and things will go better with us in every direction!

[Georg Sverdrup.]

